% !TEX program = xelatex
% =============================================================================
%  近代物理实验报告 —— LaTeX 骨架模板
%
%  用法：把本文件复制成报告工程的 main.tex，替换标题块与各节正文，
%        图放在 figs/ 下，用 xelatex 连续编译两次。
%
%  编译： xelatex -interaction=nonstopmode main.tex   （跑 2 次）
%  检查： 编译日志中不应出现 "^!"、Overfull \hbox、Reference ... undefined
%
%  版式依据《近代物理实验报告论文模板及说明》，细目见
%  references/latex-layout-spec.md；本文件编译后即是一份可直接对照的样张。
% =============================================================================
\documentclass[UTF8,zihao=5,a4paper]{ctexart}

% ----------------------------- 页面 -----------------------------
\usepackage[left=2.0cm,right=2.0cm,top=2.2cm,bottom=2.2cm]{geometry}

% ----------------------------- 宏包 -----------------------------
\usepackage{amsmath,amssymb,bm}
\usepackage{siunitx}
\usepackage{graphicx,booktabs,array,tabularx,multirow}
\usepackage{tikz,pgfplots}
\usetikzlibrary{arrows.meta,positioning}
\usepackage{fancyhdr,titlesec,enumitem,caption,float,indentfirst,needspace}
\pgfplotsset{compat=1.18}

% ------------------------- 行距与段距 ---------------------------
\linespread{1.5}
\setlength{\parindent}{2em}
\setlength{\parskip}{0pt}
\setlength{\headheight}{15pt}
\setlength{\emergencystretch}{2em}

% ----------------------------- 字体 -----------------------------
% ctexart 已内置 \heiti \songti \fangsong \kaishu，不要再 \newCJKfontfamily，
% 否则报 "Control sequence \heiti is already defined"。
\setCJKmainfont{SimSun}[BoldFont=SimHei,ItalicFont=KaiTi]
\setmainfont{Times New Roman}

% --------------------------- 页眉页脚 ---------------------------
\pagestyle{fancy}
\fancyhf{}
\fancyhead[C]{\zihao{-5}\itshape MODERN PHYSICS EXPERIMENT II}
\fancyfoot[C]{\zihao{-5}\thepage}
\renewcommand{\headrulewidth}{0.4pt}

% ----------------------------- 标题 -----------------------------
% 一级：小四黑体；二级：五号黑体。不要在 \heiti 后再加 \bfseries
% （SimHei 无粗体，会报 "Font shape ... undefined"）。
\titleformat{\section}{\zihao{-4}\heiti}{\thesection}{0.8em}{}
\titleformat{\subsection}{\zihao{5}\heiti}{\thesubsection}{0.8em}{}
\titlespacing*{\section}{0pt}{2.2ex plus .4ex minus .2ex}{1.2ex plus .2ex}
\titlespacing*{\subsection}{0pt}{1.6ex plus .3ex minus .2ex}{0.9ex plus .2ex}
% 思考题小问用 \paragraph，改为独占一行的小标题
\titleformat{\paragraph}[block]{\zihao{5}\heiti}{}{0pt}{}
\titlespacing*{\paragraph}{0pt}{1.4ex plus .3ex minus .2ex}{0.5ex plus .1ex}

% --------------------------- 图题表题 ---------------------------
% 注意：不能在 \captionsetup 里直接写 \zihao{-5}（会报 "\zihao {-5} undefined"），
% 必须先 \DeclareCaptionFont 定义一个字体命令。
\DeclareCaptionFont{xiaowu}{\zihao{-5}}
\captionsetup{font=xiaowu,labelsep=quad,justification=centering}

% --------------------------- 单位格式 ---------------------------
\sisetup{per-mode=symbol,range-phrase=--}

\begin{document}

% ============================= 标题块 =============================
\begin{center}
  {\zihao{3}\heiti 实验名称：副标题或被测物理量}\par
  \vspace{0.8em}
  {\zihao{5} 作者姓名}\par
  \vspace{0.35em}
  {\zihao{5}（化学学院，南开大学，天津 300071）}\par
  \vspace{0.35em}
  {\zihao{5} E-mail：someone@example.com}
\end{center}

\vspace{0.6em}

% =========================== 摘要与关键词 =========================
\noindent\textbf{摘要：}摘要用一段式，写清实验对象、方法、主要定量结果（含数值与单位）、结论与讨论要点，不出现“本文将”“众所周知”一类空话，也不写“由于条件所限”一类自我削弱句。一般在 200--300 字。\par
\noindent\textbf{关键词：}关键词一；关键词二；关键词三；关键词四

% ============================= 正文 =============================
\section{实验目的}
用一段或若干条说明本实验要掌握的方法、要测量的物理量与要验证的规律。避免复述实验步骤。

\section{实验原理}
\subsection{物理模型与基本关系}
先给出物理模型，再给出可测量的关系式，并说明每个符号、单位与适用条件。示例：
\begin{equation}
h\nu = g\mu_B B_0,
\qquad
g=\frac{h\nu}{\mu_B B_0}
\label{eq:example}
\end{equation}
式中 $\nu$ 为微波频率（Hz），$B_0$ 为磁感应强度（T），$\mu_B$ 为玻尔磁子，
$h$ 为普朗克常量。由式~\eqref{eq:example} 可见 $g\propto\nu/B_0$，故频率与磁场的
测量误差按相同权重进入结果。

\subsection{工作原理框图}
原理框图用 TikZ 手绘，不要直接截取手册图片；图中端子名与正文一致。
\begin{figure}[htbp]
\centering
\begin{tikzpicture}[x=1cm,y=1cm,
  b/.style={draw,align=center,minimum height=0.65cm,minimum width=1.5cm,font=\zihao{6}},
  >=Stealth]
  \node[b] (s) at (0,0)   {信号源};
  \node[b] (m) at (3.0,0) {被测样品};
  \node[b] (d) at (6.0,0) {探测器};
  \node[b] (r) at (9.0,0) {记录仪};
  \draw[->] (s)--(m);
  \draw[->] (m)--(d);
  \draw[->] (d)--(r);
  \node[font=\zihao{6},above] at (1.5,0.42) {\SI{468}{nm}};
\end{tikzpicture}
\caption{测量系统工作原理框图}
\label{fig:block}
\end{figure}

\section{实验仪器与试剂}
列仪器型号、关键指标与试剂，凡手册未给出的参数不臆造。需要时用三线表：
\begin{table}[H]
\centering
\caption{主要仪器及其作用}
\label{tab:devices}
\begin{tabularx}{\linewidth}{>{\raggedright\arraybackslash}p{3.6cm}X}
\toprule
部件 & 作用\\
\midrule
光源 & 提供激发光或探测光，写明波长与功率。\\
样品 & 写明浓度、溶剂、配制与保存条件。\\
探测器 & 写明响应波段与量程。\\
\bottomrule
\end{tabularx}
\end{table}

\section{实验内容与步骤}
按任务分节（任务一、任务二……），步骤写“怎么做”与“看什么判据”，不写“认真观察”一类空话。
软件操作写清菜单路径与参数设置，便于复现。

\section{实验结果与数据处理}
\subsection{数据与拟合}
实测点与拟合曲线分开标注；拟合函数来自物理模型，不用目测拼凑。
拟合参数、决定系数与残差要给出，必要时说明拟合区间。
\begin{figure}[htbp]
\centering
\begin{tikzpicture}
\begin{axis}[width=.80\linewidth,height=6.4cm,
  xlabel={自变量 $x$ /单位}, ylabel={因变量 $y$ /单位},
  xmin=0,xmax=26,ymin=0,ymax=360,grid=both,legend pos=south east]
\addplot[only marks,mark=*,mark size=2.4pt]
  coordinates {(5,131)(10,224)(15,288)(20,311)(25,321)};
\addlegendentry{实测}
\addplot[smooth,thick,dashed]
  coordinates {(5,131)(10,224)(15,288)(20,311)(25,321)};
\addlegendentry{平滑连线}
\end{axis}
\end{tikzpicture}
\caption{因变量随自变量的变化（虚线为平滑连接实测点的曲线）}
\label{fig:data}
\end{figure}

\subsection{结果与讨论}
把数值结果写进句子，并与理论预期比较；偏离要给出机制，且区分“可能导致”与“已证明”。

\section{误差分析与讨论}
按测量链逐项讨论：标定、分辨率、对准、样品制备、温度与环境、漂移、背景、读数与拟合方式。
有重复测量或明确仪器误差时再给不确定度；没有依据时不编造误差棒。

\section{思考题}
\paragraph{（1）第一问的题目？}
先直接回答，再给模型与实测数据的支撑。保留题目原有编号。

\paragraph{（2）第二问的题目？}
同上。

\section{实验总结与展望}
结论与数据一致，不夸大；展望只写与本实验直接相关的方向。

% ============================= 参考文献 =========================
% 只保留实际引用过的条目；用户要求删除某几条时，编号随之重排，
% 正文角标同步改。参考文献用小五号。
\needspace{11\baselineskip}
\section{参考文献}
{\zihao{-5}
\begin{enumerate}[label={[\arabic*]},leftmargin=2.2em,itemsep=0.25em]
\item 作者. 书名[M]. 版次. 出版地: 出版社, 年份.
\item 作者. 篇名[J]. 刊名, 年份, 卷(期): 起止页.
\end{enumerate}
\par}

\end{document}
